\documentclass[11pt,a4paper]{article}
\usepackage[margin=1in]{geometry}
\usepackage{mathptmx}
\usepackage{amsmath}
\usepackage{amssymb}
\usepackage{amsthm}
\usepackage{braket}
\usepackage{tikz}
\usepackage{quantikz}
\usepackage{enumitem}
\usepackage{hyperref}
\usepackage{xcolor}

\definecolor{circuitblue}{RGB}{40,80,160}
\hypersetup{colorlinks=true,linkcolor=circuitblue,urlcolor=circuitblue,citecolor=circuitblue}

\theoremstyle{definition}
\newtheorem{definition}{Definition}[section]
\newtheorem{example}{Example}[section]

\theoremstyle{remark}
\newtheorem*{intuition}{Intuition}
\newtheorem*{keypoint}{Key Point}

\newcounter{exercise}[section]
\newenvironment{exercise}{\refstepcounter{exercise}\par\medskip\noindent\textbf{Exercise \theexercise.}\rmfamily}{\medskip}

\pagestyle{plain}
\setlength{\parindent}{0pt}
\setlength{\parskip}{6pt}

\begin{document}

{\LARGE \textbf{Lab 3: Noise in Quantum Computers}}\\[8pt]
{\large QML}\\[12pt]
\hrule

\vspace{16pt}

\section{introduction}

real quantum computers r noisy. like, really noisy. every gate u apply has some probability of error, qubits lose their quantum state over time (decoherence), and measurements r not perfectly accurate. this matters a lot for quantum ML bcs our variational circuits will need to run on real hardware eventually.

in this lab u'll learn to:
\begin{itemize}[leftmargin=*]
  \item simulate different types of noise using qiskit Aer
  \item compare ideal vs noisy circuit execution
  \item understand the main error channels (bit flip, phase flip, depolarizing, amplitude damping)
  \item implement basic error mitigation techniques (ZNE and measurement error mitigation)
\end{itemize}

\begin{keypoint}
we r not doing full quantum error \textit{correction} here (that requires way more qubits). instead, we're doing error \textit{mitigation} --- clever classical post-processing tricks that reduce the impact of noise without adding extra qubits.
\end{keypoint}

\section{noise channels}

\subsection{types of noise}

the main noise channels u need to know:

\begin{definition}
the \textbf{bit flip channel} flips $\ket{0} \leftrightarrow \ket{1}$ with probability $p$. it's described by:
\[
\mathcal{E}(\rho) = (1-p)\rho + p X\rho X
\]
\end{definition}

\begin{definition}
the \textbf{phase flip channel} (dephasing) applies a Z error with probability $p$:
\[
\mathcal{E}(\rho) = (1-p)\rho + p Z\rho Z
\]
\end{definition}

\begin{definition}
the \textbf{depolarizing channel} replaces the state with the maximally mixed state with probability $p$:
\[
\mathcal{E}(\rho) = (1-p)\rho + \frac{p}{3}(X\rho X + Y\rho Y + Z\rho Z)
\]
this is the ``everything goes wrong equally'' model and it's the most commonly used for benchmarking.
\end{definition}

\begin{definition}
the \textbf{amplitude damping channel} models energy loss (like a qubit decaying from $\ket{1}$ to $\ket{0}$). it's described by Kraus operators:
\[
K_0 = \begin{pmatrix} 1 & 0 \\ 0 & \sqrt{1-\gamma} \end{pmatrix}, \quad
K_1 = \begin{pmatrix} 0 & \sqrt{\gamma} \\ 0 & 0 \end{pmatrix}
\]
where $\gamma$ is the damping probability.
\end{definition}

\begin{intuition}
bit flip = classical error (like a noisy communication channel). phase flip = purely quantum error (no classical analogue). depolarizing = uniform random error. amplitude damping = energy dissipation (most physically realistic). in practice, real devices have a mix of all of these.
\end{intuition}

\begin{exercise}
\textbf{visualizing noise channels.}
\begin{enumerate}[leftmargin=*]
  \item prepare the state $\ket{+} = H\ket{0}$ and apply each noise channel (bit flip, phase flip, depolarizing) with $p = 0.1$.
  \item for each channel, compute the output density matrix using the qiskit Aer density matrix simulator.
  \item use qiskit's \texttt{plot\_bloch\_multivector} to visualize the state before and after noise.
  \item which noise channel has the biggest effect on the $\ket{+}$ state? why does this make sense?

  to add noise in qiskit:
  \begin{itemize}[leftmargin=*]
    \item import \texttt{NoiseModel} from \texttt{qiskit\_aer.noise}
    \item import the error types: \texttt{depolarizing\_error}, \texttt{pauli\_error}, etc.
    \item create a noise model and add errors to specific gates
    \item pass the noise model to the Aer simulator when u run
  \end{itemize}
\end{enumerate}
\end{exercise}

\section{building noise models}

\subsection{custom noise models in qiskit}

here's the workflow for adding noise:
\begin{enumerate}[leftmargin=*]
  \item create a \texttt{NoiseModel()} object
  \item create error objects (e.g., \texttt{depolarizing\_error(p, num\_qubits)})
  \item add them to the noise model: \texttt{noise\_model.add\_all\_qubit\_quantum\_error(error, ['gate\_name'])}
  \item run ur circuit with the noisy simulator
\end{enumerate}

u can add different error rates to different gates. typically 2-qubit gates (like CNOT) have much higher error rates than single-qubit gates. on real IBM devices, single-qubit error rates r around $10^{-3}$ to $10^{-4}$ and CNOT error rates r around $10^{-2}$.

\begin{exercise}
\textbf{noisy bell state.}
\begin{enumerate}[leftmargin=*]
  \item build the bell state circuit from lab 1 (H + CNOT + measure)
  \item run it with no noise (ideal). record the counts.
  \item create a noise model with:
    \begin{itemize}
      \item depolarizing error $p = 0.01$ on single-qubit gates
      \item depolarizing error $p = 0.05$ on CNOT gates
    \end{itemize}
  \item run the same circuit with this noise model. record the counts.
  \item compare: in the ideal case u should get only 00 and 11. wt extra outcomes appear with noise? wt r their probabilities?
  \item increase the CNOT error rate to $p = 0.1, 0.2, 0.3$. plot the probability of getting the ``wrong'' outcomes (01 or 10) as a function of the error rate.
\end{enumerate}
\end{exercise}

\begin{exercise}
\textbf{noisy GHZ state.}
\begin{enumerate}[leftmargin=*]
  \item build GHZ circuits for $n = 3, 4, 5, 6, 7$ qubits
  \item run each one with the realistic noise model (single-qubit error $= 0.001$, CNOT error $= 0.01$)
  \item for each $n$, compute the fidelity of the noisy output with the ideal GHZ state
  \item plot fidelity vs $n$. u should see it drop as $n$ increases.
  \item explain why more qubits $=$ more noise. roughly how does fidelity scale with the number of CNOT gates?
\end{enumerate}

the circuit for 4-qubit GHZ:

\begin{center}
\begin{quantikz}
\lstick{$q_0$} & \gate{H} & \ctrl{1} & \qw & \qw & \meter{} \\
\lstick{$q_1$} & \qw & \targ{} & \ctrl{1} & \qw & \meter{} \\
\lstick{$q_2$} & \qw & \qw & \targ{} & \ctrl{1} & \meter{} \\
\lstick{$q_3$} & \qw & \qw & \qw & \targ{} & \meter{}
\end{quantikz}
\end{center}
\end{exercise}

\section{amplitude damping and $T_1$ decay}

amplitude damping is the most physically motivated noise model. it represents the qubit losing energy over time ($T_1$ relaxation).

the probability of decay depends on the gate duration $t$ and the $T_1$ time:
\[
\gamma = 1 - e^{-t/T_1}
\]

typical IBM quantum devices have $T_1 \approx 100\,\mu\text{s}$ and single-qubit gate times of $\approx 35\,\text{ns}$, so the error per gate is small. but circuits with many gates accumulate error.

\begin{exercise}
\textbf{$T_1$ decay simulation.}
\begin{enumerate}[leftmargin=*]
  \item prepare $\ket{1}$ (apply X to $\ket{0}$)
  \item add $k$ identity gates (each representing some time duration)
  \item add amplitude damping noise with $\gamma = 0.01$ per gate
  \item measure and compute $P(\ket{1})$
  \item repeat for $k = 0, 10, 20, 50, 100, 200$
  \item plot $P(\ket{1})$ vs $k$. u should see exponential decay.
  \item fit an exponential $P = e^{-k/k_0}$ and extract $k_0$. compare with $1/\gamma$.
\end{enumerate}

this is essentially how $T_1$ is measured on real devices!
\end{exercise}

\section{zero-noise extrapolation (ZNE)}

ZNE is one of the simplest and most effective error mitigation techniques. the idea is:
\begin{enumerate}[leftmargin=*]
  \item run ur circuit at several different noise levels
  \item extrapolate back to zero noise
\end{enumerate}

but wait --- u can't just turn down the noise on a real device. the trick is to \textit{increase} the noise by stretching the circuit (adding more gates that don't change the ideal computation but do add noise).

\subsection{noise amplification by unitary folding}

the simplest way to amplify noise: replace each gate $G$ with $G \cdot G^\dagger \cdot G$. this does the same thing as $G$ (bcs $G^\dagger G = I$) but uses 3x the gates, so it has 3x the noise. u can generalize: $G \to G(G^\dagger G)^k$ gives a noise scale factor of $2k+1$.

here's wt unitary folding looks like for a simple circuit:

\begin{center}
original circuit:

\begin{quantikz}
\lstick{$q_0$} & \gate{H} & \ctrl{1} & \meter{} \\
\lstick{$q_1$} & \qw & \targ{} & \meter{}
\end{quantikz}

\vspace{8pt}
folded once ($3\times$ noise):

\begin{quantikz}
\lstick{$q_0$} & \gate{H} & \ctrl{1} & \ctrl{1} & \gate{H} & \gate{H} & \ctrl{1} & \meter{} \\
\lstick{$q_1$} & \qw & \targ{} & \targ{} & \qw & \qw & \targ{} & \meter{}
\end{quantikz}
\end{center}

\begin{exercise}
\textbf{implementing ZNE.}
\begin{enumerate}[leftmargin=*]
  \item build a simple circuit: prepare $\ket{+}$ and measure the expectation value of $Z$ (should be 0 ideally).
  \item add depolarizing noise with $p = 0.05$ on all gates.
  \item run the circuit normally (noise factor $= 1$). record $\braket{Z}$.
  \item fold the circuit once (noise factor $= 3$): for each gate $G$, replace it with $G \cdot G^\dagger \cdot G$. record $\braket{Z}$.
  \item fold twice (noise factor $= 5$). record $\braket{Z}$.
  \item now extrapolate to noise factor $= 0$:
    \begin{itemize}
      \item u have data points at noise factors 1, 3, 5
      \item fit a line (or quadratic) through these points
      \item evaluate the fit at noise factor $= 0$
    \end{itemize}
  \item compare ur extrapolated value with the ideal value. how much closer r u?
\end{enumerate}

hint for folding in qiskit: create a new circuit. for each gate in the original circuit, add the gate, then its inverse, then the gate again. the inverse of H is H. the inverse of CNOT is CNOT. the inverse of $R_y(\theta)$ is $R_y(-\theta)$.
\end{exercise}

\begin{exercise}
\textbf{ZNE on a bell state.}
apply ZNE to the bell state circuit:
\begin{enumerate}[leftmargin=*]
  \item measure $\braket{ZZ}$ on the bell state (should be $+1$ ideally)
  \item run at noise factors 1, 3, 5, 7
  \item extrapolate to zero noise using:
    \begin{itemize}
      \item linear fit (using noise factors 1 and 3 only)
      \item quadratic fit (using noise factors 1, 3, and 5)
    \end{itemize}
  \item which fit gives a better estimate? plot ur data points and both fits.
  \item wt happens if u try noise factors that r too large (e.g., 15, 21)? why does ZNE break down at very high noise?
\end{enumerate}
\end{exercise}

\section{measurement error mitigation}

another common source of error: the measurement itself isn't perfect. on real hardware, u might measure $\ket{0}$ when the qubit is actually in $\ket{1}$, or vice versa.

the idea behind measurement error mitigation:
\begin{enumerate}[leftmargin=*]
  \item characterize the measurement errors by preparing known states and measuring them
  \item build a \textit{calibration matrix} $M$ where $M_{ij} = P(\text{measure } i \mid \text{prepared } j)$
  \item invert this matrix to correct future measurements: $\vec{p}_{\text{corrected}} = M^{-1}\vec{p}_{\text{measured}}$
\end{enumerate}

for 1 qubit, the calibration matrix is:
\[
M = \begin{pmatrix}
P(0|0) & P(0|1) \\
P(1|0) & P(1|1)
\end{pmatrix}
\]

\begin{exercise}
\textbf{measurement error mitigation.}
\begin{enumerate}[leftmargin=*]
  \item create a noise model that only has measurement errors:
    \begin{itemize}
      \item $P(\text{measure 0} | \text{state is 1}) = 0.05$ (5\% chance of 1 being read as 0)
      \item $P(\text{measure 1} | \text{state is 0}) = 0.03$ (3\% chance of 0 being read as 1)
    \end{itemize}
  \item build the calibration matrix by:
    \begin{itemize}
      \item preparing $\ket{0}$ and measuring many times $\to$ gives column 0 of $M$
      \item preparing $\ket{1}$ and measuring many times $\to$ gives column 1 of $M$
    \end{itemize}
  \item now prepare the state $\ket{+}$ and measure. compare:
    \begin{itemize}
      \item ideal result (no noise)
      \item noisy result (with measurement error)
      \item mitigated result (apply $M^{-1}$)
    \end{itemize}
  \item extend this to 2 qubits. how big is the calibration matrix now? how many calibration circuits do u need?
\end{enumerate}
\end{exercise}

\begin{exercise}
\textbf{combining noise and mitigation.}
this one puts everything together:
\begin{enumerate}[leftmargin=*]
  \item build a 3-qubit GHZ circuit
  \item add a realistic noise model: depolarizing ($p = 0.005$ single, $p = 0.02$ CNOT) plus measurement error (3\%)
  \item run and record the raw noisy results
  \item apply measurement error mitigation
  \item apply ZNE (noise factors 1, 3, 5)
  \item compare all four results: ideal, noisy, measurement-mitigated only, ZNE + measurement-mitigated
  \item which technique helps more? does combining them help even more?
\end{enumerate}

\begin{center}
\begin{quantikz}
\lstick{$q_0$} & \gate{H} & \ctrl{1} & \qw & \qw & \meter{} \\
\lstick{$q_1$} & \qw & \targ{} & \ctrl{1} & \qw & \meter{} \\
\lstick{$q_2$} & \qw & \qw & \targ{} & \qw & \meter{}
\end{quantikz}
\end{center}
\end{exercise}

\section{noise and circuit depth}

\begin{exercise}
\textbf{depth vs fidelity.}
this exercise builds intuition for why shallow circuits r so important in quantum ML:
\begin{enumerate}[leftmargin=*]
  \item build circuits of increasing depth: create a random circuit using repeated layers of $R_y(\theta)$ on each qubit followed by a CNOT entangling layer. use 4 qubits.
  \item run each circuit with a realistic noise model and compute the output state fidelity.
  \item plot fidelity vs circuit depth (number of layers).
  \item at wt depth does fidelity drop below 0.5? below 0.9?
  \item now imagine u r designing a variational circuit for ML. based on ur findings, wt's a reasonable maximum depth?
\end{enumerate}

this is one of the key constraints in quantum ML: ur variational circuits can't be too deep or the noise will destroy all the information. keep this in mind for lab 5.
\end{exercise}

\section{submission}

submit a jupyter notebook with:
\begin{itemize}[leftmargin=*]
  \item all exercises completed
  \item plots comparing ideal vs noisy vs mitigated results
  \item the depth vs fidelity analysis
  \item a short reflection: how does noise affect the feasibility of quantum ML on near-term devices?
\end{itemize}

\end{document}
